Ziel dieser Arbeit ist die Untersuchung, ob ein automatisierter Feddbackloop zur Verbesserung der Qualität von Texten in Leichter Sprache Plus geeignet ist. Hierfür wird ein Proof of Concept konzipiert und prototypisch umgesetzt, der generierte LS+-Texte anhand definierter Qualitätsmetriken bewertet und auf Basis der Bewertung automatisierte Verbesserungen ableitet.

Ein Schwerpunkt der Arbeit liegt auf der Analyse geeigneter Bewertungsverfahren für Texte in Leichter Sprache Plus. Dabei soll untersucht werden, welche linguistischen, regelbasierten und semantischen Metriken geeignet sind, um die Qualität von LS+-Texten automatisiert zu erfassen.
Die gewonnenen Erkenntnisse dienen als Grundlage für einen Feedbackloop, der erkannte schwächen im Text an ein Large Language Model zurückmeldet und dadurch eine iterative Optimierung der Texte ermöglicht.

Im Rahmen der Arbeit wird nicht die Entwicklung eines produktionsreifen Systems angestrebt.
Vielmehr soll untersucht werden, ob  ein solcher Ansatz grundsätzlich geeignet ist, die Qualität automatisch erzeugter LS+-Texte zu verbessern und welche Möglichkeiten sowie Grenzen sich dabei ergeben.

Aus der Zielsetzung ergeben sich folgende Forschungsfragen:
\begin{itemize}
	\item \textbf{F1:} Welche automatisiert erfassbaren Metriken eignen sich zur Bewertung der Qualität von Texten in Leichter Sprache Plus?
	
	\item \textbf{F2:} Inwieweit kann ein auf Metriken basierter Feedbackloop die Qualität automatisch erzeugter LS+-Texte verbessern und sicherstellen?
	
	\item \textbf{F3:} Welche Auswirkungen hat die iterative Überarbeitung auf Verständlichkeit, Regelkonformität und inhaltliche Konsistenz der erzeugten Texte?
\end{itemize}

Zur Beantwortung dieser Forschungsfragen wird ein Proof of Concept entwickelt und anhand verschiedener Textbeispiele ausgewertet.
Dabei werden die Ergebnisse vor und nach Anwendung des Feedbackloops miteinander verglichen.
Zusätzlich erfolgt ein Vergleich mit bereits vorhandenen LS+-Texten, um die Wirksamkeit und tatsächliche Nutzbarkeit einschätzen zu können.